
% Learning objectives
\begin{frame}{Learning Objectives}
By the end of this chapter, you will understand:
\begin{itemize}
\item Instruction Level Parallelism (ILP) concepts and benefits
\item Pipelining principles and implementation techniques
\item Pipeline stages and timing analysis
\item Hazards: structural, data, and control dependencies
\item Advanced ILP techniques: superscalar, out-of-order execution
\item Performance metrics and optimization strategies
\item Trade-offs between complexity and performance
\end{itemize}
\end{frame}

% subsection 1: Introduction to ILP
\subsection{Introduction to ILP}

\begin{frame}{What is Instruction Level Parallelism?}
\begin{columns}
\begin{column}{0.6\textwidth}
\concept{ILP Definition}:
\begin{itemize}
\item Parallel execution of instructions within a single program
\item Exploiting independence between instructions
\item Multiple instructions in various stages of execution
\item Fundamental technique for performance improvement
\end{itemize}

\vspace{0.3cm}
\concept{Forms of ILP}:
\begin{itemize}
\item \textbf{Pipelining}: Overlapping instruction execution stages
\item \textbf{Multiple Issue}: Multiple instructions per cycle
\item \textbf{Out-of-Order Execution}: Dynamic instruction reordering
\item \textbf{Speculative Execution}: Predicting and executing ahead
\end{itemize}

\vspace{0.3cm}
\concept{Performance Benefits}:
\begin{itemize}
\item Higher instruction throughput
\item Better resource utilization
\item Improved clock frequency potential
\item Enhanced overall system performance
\end{itemize}

\vspace{0.3cm}
\concept{Implementation Challenges}:
\begin{itemize}
\item Managing dependencies between instructions
\item Handling control flow changes
\item Increased hardware complexity
\item Power consumption considerations
\end{itemize}
\end{column}
\begin{column}{0.4\textwidth}
\begin{center}
\textbf{ILP Techniques}
\vspace{0.3cm}

\tikz{
% Central ILP concept
\node[draw, circle, minimum width=1.5cm] (ilp) at (0,0) {ILP};

% Techniques around it
\node[draw, rectangle] (pipe) at (0,1.8) {Pipelining};
\node[draw, rectangle] (multi) at (-1.8,0.9) {Multiple};
\node[draw, rectangle] (multi2) at (-1.8,0.6) {Issue};
\node[draw, rectangle] (ooo) at (1.8,0.9) {Out-of-Order};
\node[draw, rectangle] (ooo2) at (1.8,0.6) {Execution};
\node[draw, rectangle] (spec) at (0,-1.8) {Speculative};
\node[draw, rectangle] (spec2) at (0,-2.1) {Execution};

% Connections
\draw[->] (pipe) -- (ilp);
\draw[->] (multi) -- (ilp);
\draw[->] (ooo) -- (ilp);
\draw[->] (spec) -- (ilp);
}
\end{center}
\end{column}
\end{columns}
\end{frame}

\begin{frame}{Sequential vs. Parallel Execution}
\begin{columns}
\begin{column}{0.6\textwidth}
\concept{Sequential Execution}:
\begin{itemize}
\item One instruction completes before next begins
\item Simple control logic and implementation
\item Underutilized hardware resources
\item Lower performance potential
\end{itemize}

\vspace{0.3cm}
\concept{Parallel Execution (ILP)}:
\begin{itemize}
\item Multiple instructions in flight simultaneously
\item Complex control and dependency management
\item Better hardware resource utilization
\item Higher performance potential
\end{itemize}

\vspace{0.3cm}
\concept{Key Differences}:
\begin{itemize}
\item \textbf{Throughput}: Instructions per cycle
\item \textbf{Latency}: Time for single instruction
\item \textbf{Resource Usage}: Functional unit utilization
\item \textbf{Complexity}: Control logic requirements
\end{itemize}

\vspace{0.3cm}
\concept{Performance Metrics}:
\begin{itemize}
\item \textbf{CPI}: Cycles per instruction
\item \textbf{IPC}: Instructions per cycle
\item \textbf{Throughput}: Instructions per second
\item \textbf{Speedup}: Performance improvement ratio
\end{itemize}
\end{column}
\begin{column}{0.4\textwidth}
\begin{center}
\textbf{Execution Comparison}
\vspace{0.3cm}

% Sequential execution
\tikz{
\node at (1.5,3.5) {\textbf{Sequential}};

% Instructions
\foreach \x/\inst in {0/I1, 0.8/I2, 1.6/I3, 2.4/I4} {
  \draw[thick, fill=blue!30] (\x,3) rectangle (\x+0.6,3.3);
  \node at (\x+0.3,3.15) {\inst};
}

% Time arrow
\draw[->] (0,2.7) -- (3,2.7);
\node at (1.5,2.5) {Time};
\node at (1.5,2.2) {4 cycles};
}

\vspace{0.5cm}

% Parallel execution
\tikz{
\node at (1.5,2.5) {\textbf{Pipelined}};

% Pipeline stages
\foreach \y/\stage in {2.2/IF, 1.9/ID, 1.6/EX, 1.3/WB} {
  \node at (-0.3,\y) {\stage};
}

% Instructions in pipeline
\foreach \x/\cycle in {0/1, 0.4/2, 0.8/3, 1.2/4} {
  % I1
  \draw[thick, fill=red!30] (\x,2.1) rectangle (\x+0.3,2.3);
  % I2
  \if\cycle1\else
    \draw[thick, fill=green!30] (\x,1.8) rectangle (\x+0.3,2);
  \fi
  % I3
  \if\cycle1\else\if\cycle2\else
    \draw[thick, fill=blue!30] (\x,1.5) rectangle (\x+0.3,1.7);
  \fi\fi
  % I4
  \if\cycle1\else\if\cycle2\else\if\cycle3\else
    \draw[thick, fill=yellow!30] (\x,1.2) rectangle (\x+0.3,1.4);
  \fi\fi\fi
}

% Time arrow
\draw[->] (0,0.9) -- (1.5,0.9);
\node at (0.75,0.7) {Time};
\node at (0.75,0.4) {4 cycles, 4 instructions};
}
\end{center}
\end{column}
\end{columns}
\end{frame}

% subsection 2: Pipelining Fundamentals
\subsection{Pipelining Fundamentals}

\begin{frame}{Pipeline Acceleration Principle}
\begin{columns}
\begin{column}{0.6\textwidth}
\concept{Basic Principle}:
\begin{itemize}
\item Break complex operation into simpler stages
\item Insert pipeline registers between stages
\item Allow multiple operations to overlap
\item Increase clock frequency by reducing critical path
\end{itemize}

\vspace{0.3cm}
\concept{Timing Analysis}:
\begin{itemize}
\item \textbf{Without Pipeline}: $T_{clock} = t_{CLC_{f \circ g}}$
\item \textbf{With Pipeline}: $T_{clock} = \max(t_{CLC_f}, t_{CLC_g})$
\item \textbf{Acceleration}: $\alpha = \frac{t_{CLC_{f \circ g}}}{\max(t_{CLC_f}, t_{CLC_g})}$
\item \textbf{Maximum Efficiency}: $\alpha \approx 2$ when $t_{CLC_f} \approx t_{CLC_g}$
\end{itemize}

\vspace{0.3cm}
\concept{Design Considerations}:
\begin{itemize}
\item Balance pipeline stage delays
\item Minimize register overhead
\item Consider setup and hold times
\item Account for clock skew and jitter
\end{itemize}

\vspace{0.3cm}
\concept{Scalability}:
\begin{itemize}
\item Process can be applied recursively
\item Continue until desired frequency achieved
\item Diminishing returns with more stages
\item Increased latency with deeper pipelines
\end{itemize}
\end{column}
\begin{column}{0.4\textwidth}
\begin{center}
\textbf{Pipeline Acceleration}
\vspace{0.3cm}

% Non-pipelined circuit
\tikz{
\node at (1.5,3.5) {\textbf{Non-Pipelined}};

\draw[thick] (0,3) rectangle (0.5,3.3);
\node at (0.25,3.15) {Reg};

\draw[thick] (0.7,3) rectangle (2.3,3.3);
\node at (1.5,3.15) {$CLC_{f \circ g}$};

\draw[thick] (2.5,3) rectangle (3,3.3);
\node at (2.75,3.15) {Reg};

\draw[->] (0.5,3.15) -- (0.7,3.15);
\draw[->] (2.3,3.15) -- (2.5,3.15);

\node at (1.5,2.7) {$T_{clock} = t_{CLC_{f \circ g}}$};
}

\vspace{0.5cm}

% Pipelined circuit
\tikz{
\node at (1.5,2.5) {\textbf{Pipelined}};

\draw[thick] (0,2) rectangle (0.4,2.3);
\node at (0.2,2.15) {Reg};

\draw[thick] (0.6,2) rectangle (1.2,2.3);
\node at (0.9,2.15) {$CLC_f$};

\draw[thick] (1.4,2) rectangle (1.8,2.3);
\node at (1.6,2.15) {Reg};

\draw[thick] (2,2) rectangle (2.6,2.3);
\node at (2.3,2.15) {$CLC_g$};

\draw[thick] (2.8,2) rectangle (3.2,2.3);
\node at (3,2.15) {Reg};

\draw[->] (0.4,2.15) -- (0.6,2.15);
\draw[->] (1.2,2.15) -- (1.4,2.15);
\draw[->] (1.8,2.15) -- (2,2.15);
\draw[->] (2.6,2.15) -- (2.8,2.15);

\node at (1.6,1.7) {$T_{clock} = \max(t_{CLC_f}, t_{CLC_g})$};
\node at (1.6,1.4) {Speedup: $\alpha \approx 2$};
}
\end{center}
\end{column}
\end{columns}
\end{frame}

\begin{frame}{Classic RISC Pipeline Stages}
\begin{columns}
\begin{column}{0.6\textwidth}
\concept{Five-Stage Pipeline}:
\begin{enumerate}
\item \textbf{IF (Instruction Fetch)}: Fetch instruction from memory
\item \textbf{ID (Instruction Decode)}: Decode instruction and read registers
\item \textbf{EX (Execute)}: Perform ALU operation or address calculation
\item \textbf{MEM (Memory Access)}: Access data memory if needed
\item \textbf{WB (Write Back)}: Write result back to register file
\end{enumerate}

\vspace{0.3cm}
\concept{Stage Characteristics}:
\begin{itemize}
\item Each stage takes one clock cycle
\item Pipeline registers store intermediate results
\item Control signals propagated through pipeline
\item Data forwarding may be needed between stages
\end{itemize}

\vspace{0.3cm}
\concept{Pipeline Benefits}:
\begin{itemize}
\item Up to 5x speedup in ideal conditions
\item One instruction completed per cycle (after fill)
\item Better resource utilization
\item Scalable to more complex processors
\end{itemize}

\vspace{0.3cm}
\concept{Pipeline Challenges}:
\begin{itemize}
\item Pipeline hazards and dependencies
\item Branch prediction requirements
\item Exception handling complexity
\item Increased design verification effort
\end{itemize}
\end{column}
\begin{column}{0.4\textwidth}
\begin{center}
\textbf{5-Stage Pipeline}
\vspace{0.3cm}

\tikz{
% Pipeline stages
\foreach \y/\stage/\desc in {
  3.5/IF/Fetch,
  3/ID/Decode,
  2.5/EX/Execute,
  2/MEM/Memory,
  1.5/WB/Write Back
} {
  \draw[thick] (0,\y-0.2) rectangle (2,\y+0.2);
  \node at (0.3,\y) {\textbf{\stage}};
  \node at (1.3,\y) {\desc};
}

% Pipeline registers
\foreach \y in {3.3, 2.8, 2.3, 1.8} {
  \draw[thick, fill=yellow!30] (2.2,\y-0.1) rectangle (2.6,\y+0.1);
  \node at (2.4,\y) {Reg};
}

% Data flow
\foreach \y in {3.5, 3, 2.5, 2, 1.5} {
  \draw[->] (2,\y) -- (2.2,\y);
}
\draw[->] (2.6,1.8) -- (2.6,1.3);
\draw[->] (2.6,1.3) -- (0,1.3);
\draw[->] (0,1.3) -- (0,1.7);

% Timing
\node at (1,1) {\textbf{Latency: 5 cycles}};
\node at (1,0.7) {\textbf{Throughput: 1 inst/cycle}};
}
\end{center}
\end{column}
\end{columns}
\end{frame}

% subsection 3: Pipeline Implementation
\subsection{Pipeline Implementation}

\begin{frame}{toyRISC Pipelined Implementation}
\begin{columns}
\begin{column}{0.6\textwidth}
\concept{Four-Stage toyRISC Pipeline}:
\begin{enumerate}
\item \textbf{IF (Instruction Fetch)}: PC addresses program memory, compute next PC
\item \textbf{OF (Operands Fetch)}: Read operands from register file
\item \textbf{EX (Execution)}: ALU operation, branch address calculation
\item \textbf{WB (Write Back)}: Update PC and write result to register file
\end{enumerate}

\vspace{0.3cm}
\concept{Timing Analysis}:
\begin{itemize}
\item $t_{IF} = t_{pc} + \max(t_{next}, t_{ProgMem}) + t_{PipeReg1}$
\item $t_{OF} = t_{PipeReg1} + t_{regFile} + t_{PipeReg2}$
\item $t_{EX} = t_{PipeReg2} + \max(t_{ALU}, t_{add}, t_{DataMem}) + t_{PipeReg3}$
\item $t_{WB} = t_{PipeReg3} + \max(t_{regFile}, t_{PC})$
\end{itemize}

\vspace{0.3cm}
\concept{Clock Frequency}:
\begin{itemize}
\item $T_{clock} = \max(t_{IF}, t_{OF}, t_{EX}, t_{WB})$
\item Balanced pipeline stages for optimal performance
\item Critical path determines maximum frequency
\end{itemize}
\end{column}
\begin{column}{0.4\textwidth}
\begin{center}
\textbf{toyRISC Pipeline}
\vspace{0.3cm}

\tikz{
% Pipeline stages
\node at (1.5,3.8) {\textbf{toyRISC 4-Stage}};

% IF stage
\draw[thick] (0,3.3) rectangle (3,3.6);
\node at (0.2,3.45) {\textbf{IF}};
\node at (1.5,3.45) {PC → ProgMem, NextPC};

% OF stage
\draw[thick] (0,2.8) rectangle (3,3.1);
\node at (0.2,2.95) {\textbf{OF}};
\node at (1.5,2.95) {RegFile Read};

% EX stage
\draw[thick] (0,2.3) rectangle (3,2.6);
\node at (0.2,2.45) {\textbf{EX}};
\node at (1.5,2.45) {ALU, DataMem, Branch};

% WB stage
\draw[thick] (0,1.8) rectangle (3,2.1);
\node at (0.2,1.95) {\textbf{WB}};
\node at (1.5,1.95) {RegFile Write, PC Update};

% Pipeline registers
\foreach \y in {3.15, 2.65, 2.15} {
  \draw[thick, fill=yellow!30] (3.2,\y-0.1) rectangle (3.6,\y+0.1);
  \node at (3.4,\y) {Reg};
}

% Latency
\node at (1.5,1.4) {\textbf{Latency: 3 cycles}};
\node at (1.5,1.1) {\textbf{4 instructions in flight}};
}
\end{center}
\end{column}
\end{columns}
\end{frame}

\begin{frame}[fragile]{Pipeline Execution Example}
\begin{columns}
\begin{column}{0.6\textwidth}
\concept{Instruction Sequence}:
%\begin{lstlisting}[style=Assembler]
\begin{lstlisting}
XOR(5,0,0)  ; R5 = R0 XOR R0
VAL(1,12)   ; R1 = 12
SUB(2,3,4)  ; R2 = R3 - R4
ADD(0,0,3)  ; R0 = R0 + R3
AND(1,3,2)  ; R1 = R3 AND R2
\end{lstlisting}

\vspace{0.3cm}
\concept{Pipeline Execution}:
\begin{itemize}
\item Each instruction takes 4 cycles to complete
\item New instruction starts every cycle
\item After initial fill, one instruction completes per cycle
\item Steady-state throughput: 1 instruction/cycle
\end{itemize}

\vspace{0.3cm}
\concept{Performance Analysis}:
\begin{itemize}
\item \textbf{Latency}: 3 cycles (pipeline depth - 1)
\item \textbf{Throughput}: 1 instruction per cycle
\item \textbf{Speedup}: Approximately 4x over non-pipelined
\item \textbf{Efficiency}: Depends on hazard frequency
\end{itemize}
\end{column}
\begin{column}{0.4\textwidth}
\begin{center}
\textbf{Pipeline Timeline}
\vspace{0.3cm}

\begin{table}[h]
\tiny
\begin{tabular}{|c||c|c|c|c|c|c|c|}
\hline
\textbf{Stage/Time} & \textbf{i} & \textbf{i+1} & \textbf{i+2} & \textbf{i+3} & \textbf{i+4} & \textbf{i+5} & \textbf{i+6} \\
\hline\hline
\textbf{IF} & XOR & VAL & SUB & ADD & AND & & \\
\hline
\textbf{OF} & & XOR & VAL & SUB & ADD & AND & \\
\hline
\textbf{EX} & & & XOR & VAL & SUB & ADD & AND \\
\hline
\textbf{WB} & & & & XOR & VAL & SUB & ADD \\
\hline
\end{tabular}
\end{table}

\vspace{0.3cm}

\tikz{
% Completion timeline
\node at (1.5,2.5) {\textbf{Instruction Completion}};

\foreach \x/\inst in {0/XOR, 0.6/VAL, 1.2/SUB, 1.8/ADD, 2.4/AND} {
  \draw[thick, fill=green!30] (\x,2) rectangle (\x+0.5,2.3);
  \node at (\x+0.25,2.15) {\inst};
}

\draw[->] (0,1.7) -- (2.9,1.7);
\node at (1.45,1.5) {Time (cycles)};

\node at (1.5,1.2) {\textbf{Steady State:}};
\node at (1.5,0.9) {1 instruction/cycle};
}
\end{center}
\end{column}
\end{columns}
\end{frame}

% subsection 4: Hazards and Dependencies
\subsection{Hazards and Dependencies}

\begin{frame}{Types of Pipeline Hazards}
\begin{columns}
\begin{column}{0.6\textwidth}
\concept{Structural Hazards}:
\begin{itemize}
\item Hardware resource conflicts
\item Multiple instructions need same resource
\item Examples: memory port conflicts, functional unit conflicts
\item Solutions: duplicate resources, pipeline stalls
\end{itemize}

\vspace{0.3cm}
\concept{Data Hazards}:
\begin{itemize}
\item Dependencies between instructions
\item \textbf{RAW (Read After Write)}: True dependency
\item \textbf{WAR (Write After Read)}: Anti-dependency
\item \textbf{WAW (Write After Write)}: Output dependency
\end{itemize}

\vspace{0.3cm}
\concept{Control Hazards}:
\begin{itemize}
\item Branch and jump instructions
\item Unknown next instruction address
\item Pipeline must stall or speculate
\item Solutions: branch prediction, delayed branches
\end{itemize}

\vspace{0.3cm}
\concept{Hazard Resolution}:
\begin{itemize}
\item \textbf{Stalling}: Insert pipeline bubbles
\item \textbf{Forwarding}: Bypass pipeline registers
\item \textbf{Prediction}: Speculative execution
\item \textbf{Reordering}: Dynamic instruction scheduling
\end{itemize}
\end{column}
\begin{column}{0.4\textwidth}
\begin{center}
\textbf{Hazard Types}
\vspace{0.3cm}

% Structural hazard
\tikz{
\node at (1.5,3.5) {\textbf{Structural}};
\draw[thick] (0,3.1) rectangle (3,3.3);
\node at (1.5,3.2) {Inst1: MEM access};
\draw[thick] (0,2.9) rectangle (3,3.1);
\node at (1.5,3) {Inst2: MEM access};
\draw[thick, red] (2.5,2.9) rectangle (2.7,3.3);
\node at (2.6,2.7) {Conflict};
}

\vspace{0.3cm}

% Data hazard
\tikz{
\node at (1.5,2.5) {\textbf{Data (RAW)}};
\draw[thick] (0,2.1) rectangle (3,2.3);
\node at (1.5,2.2) {ADD R1, R2, R3};
\draw[thick] (0,1.9) rectangle (3,2.1);
\node at (1.5,2) {SUB R4, R1, R5};
\draw[->, red, thick] (0.8,2.1) -- (1.2,1.9);
\node at (2.6,1.7) {R1 dep};
}

\vspace{0.3cm}

% Control hazard
\tikz{
\node at (1.5,1.5) {\textbf{Control}};
\draw[thick] (0,1.1) rectangle (3,1.3);
\node at (1.5,1.2) {BEQ R1, R2, Label};
\draw[thick] (0,0.9) rectangle (3,1.1);
\node at (1.5,1) {??? (unknown)};
\draw[->, red, thick] (1.5,1.1) -- (1.5,0.9);
\node at (2.6,0.7) {Branch};
}
\end{center}
\end{column}
\end{columns}
\end{frame}

\begin{frame}{Data Dependencies and Forwarding}
\begin{columns}
\begin{column}{0.6\textwidth}
\concept{Read After Write (RAW) Dependency}:
\begin{itemize}
\item Most common and critical dependency
\item Instruction reads register written by previous instruction
\item Creates pipeline stall without forwarding
\item Example: ADD R1, R2, R3; SUB R4, R1, R5
\end{itemize}

\vspace{0.3cm}
\concept{Forwarding/Bypassing Solution}:
\begin{itemize}
\item Forward result from later pipeline stage
\item Bypass intermediate pipeline registers
\item Reduces or eliminates pipeline stalls
\item Requires additional hardware paths
\end{itemize}

\vspace{0.3cm}
\concept{Forwarding Paths}:
\begin{itemize}
\item \textbf{EX-to-EX}: ALU result to ALU input
\item \textbf{MEM-to-EX}: Memory result to ALU input
\item \textbf{WB-to-EX}: Write-back result to ALU input
\item \textbf{Load-Use}: Special case requiring stall
\end{itemize}

\vspace{0.3cm}
\concept{Limitations}:
\begin{itemize}
\item Load-use hazards still require stalls
\item Increased hardware complexity
\item Additional multiplexers and control logic
\item Power consumption overhead
\end{itemize}
\end{column}
\begin{column}{0.4\textwidth}
\begin{center}
\textbf{Data Forwarding}
\vspace{0.3cm}

% Without forwarding
\tikz{
\node at (1.5,3.5) {\textbf{Without Forwarding}};

% Pipeline stages
\foreach \y/\stage in {3.2/EX, 2.9/MEM, 2.6/WB} {
  \draw[thick] (0,\y-0.1) rectangle (1,\y+0.1);
  \node at (0.5,\y) {\stage};
}

% Instructions
\draw[thick, fill=red!30] (1.2,3.1) rectangle (1.8,3.3);
\node at (1.5,3.2) {ADD};
\draw[thick, fill=blue!30] (1.2,2.8) rectangle (1.8,3);
\node at (1.5,2.9) {SUB};

% Stall
\draw[thick, fill=gray!30] (2,2.8) rectangle (2.6,3);
\node at (2.3,2.9) {STALL};

\node at (1.5,2.3) {2 cycle stall};
}

\vspace{0.5cm}

% With forwarding
\tikz{
\node at (1.5,2) {\textbf{With Forwarding}};

% Pipeline stages
\foreach \y/\stage in {1.7/EX, 1.4/MEM, 1.1/WB} {
  \draw[thick] (0,\y-0.1) rectangle (1,\y+0.1);
  \node at (0.5,\y) {\stage};
}

% Instructions
\draw[thick, fill=red!30] (1.2,1.6) rectangle (1.8,1.8);
\node at (1.5,1.7) {ADD};
\draw[thick, fill=blue!30] (1.2,1.3) rectangle (1.8,1.5);
\node at (1.5,1.4) {SUB};

% Forwarding path
\draw[->, green, thick] (1.8,1.7) -- (2.2,1.7) -- (2.2,1.4) -- (1.8,1.4);
\node at (2.5,1.55) {Forward};

\node at (1.5,0.8) {No stall};
}
\end{center}
\end{column}
\end{columns}
\end{frame}

\begin{frame}{Control Hazards and Branch Prediction}
\begin{columns}
\begin{column}{0.6\textwidth}
\concept{Control Hazard Problem}:
\begin{itemize}
\item Branch outcome unknown until EX stage
\item Next instruction address uncertain
\item Pipeline must stall or guess
\item Significant performance impact
\end{itemize}

\vspace{0.3cm}
\concept{Branch Prediction Techniques}:
\begin{itemize}
\item \textbf{Static Prediction}: Always taken/not taken
\item \textbf{Dynamic Prediction}: Based on history
\item \textbf{1-bit Predictor}: Last outcome
\item \textbf{2-bit Predictor}: Saturating counter
\item \textbf{Correlating Predictors}: Pattern-based
\end{itemize}

\vspace{0.3cm}
\concept{Branch Target Prediction}:
\begin{itemize}
\item \textbf{Branch Target Buffer (BTB)}: Cache of target addresses
\item \textbf{Return Address Stack}: For function calls
\item \textbf{Indirect Branch Prediction}: For computed jumps
\end{itemize}

\vspace{0.3cm}
\concept{Misprediction Recovery}:
\begin{itemize}
\item Flush incorrectly fetched instructions
\item Restore processor state
\item Restart from correct address
\item Performance penalty for mispredictions
\end{itemize}
\end{column}
\begin{column}{0.4\textwidth}
\begin{center}
\textbf{Branch Prediction}
\vspace{0.3cm}

% 2-bit predictor state machine
\tikz{
\node at (1.5,3.5) {\textbf{2-bit Predictor}};

% States
\node[draw, circle] (st) at (0,2.5) {ST};
\node[draw, circle] (wt) at (3,2.5) {WT};
\node[draw, circle] (wn) at (0,1.5) {WN};
\node[draw, circle] (sn) at (3,1.5) {SN};

% Transitions
\draw[->] (st) to[bend left] node[above] {T} (wt);
\draw[->] (wt) to[bend left] node[below] {NT} (st);
\draw[->] (wt) to[bend left] node[right] {T} (sn);
\draw[->] (sn) to[bend left] node[left] {NT} (wt);
\draw[->] (sn) to[bend left] node[below] {T} (wn);
\draw[->] (wn) to[bend left] node[above] {NT} (sn);
\draw[->] (wn) to[bend left] node[left] {T} (st);
\draw[->] (st) to[bend left] node[right] {NT} (wn);

% Labels
\node at (1.5,0.8) {ST: Strong Taken};
\node at (1.5,0.5) {WT: Weak Taken};
\node at (1.5,0.2) {WN: Weak Not Taken};
\node at (1.5,-0.1) {SN: Strong Not Taken};
}
\end{center}
\end{column}
\end{columns}
\end{frame}

% subsection 5: Advanced ILP Techniques
\subsection{Advanced ILP Techniques}

\begin{frame}{Superscalar Processors}
\begin{columns}
\begin{column}{0.6\textwidth}
\concept{Superscalar Concept}:
\begin{itemize}
\item Multiple instructions issued per cycle
\item Multiple functional units
\item Dynamic instruction scheduling
\item Higher instruction-level parallelism
\end{itemize}

\vspace{0.3cm}
\concept{Implementation Challenges}:
\begin{itemize}
\item \textbf{Instruction Fetch}: Multiple instructions per cycle
\item \textbf{Decode Complexity}: Parallel instruction decode
\item \textbf{Resource Conflicts}: Managing multiple functional units
\item \textbf{Register Renaming}: Eliminate false dependencies
\end{itemize}

\vspace{0.3cm}
\concept{Scheduling Approaches}:
\begin{itemize}
\item \textbf{Static Scheduling}: Compiler-based reordering
\item \textbf{Dynamic Scheduling}: Hardware-based reordering
\item \textbf{Out-of-Order Execution}: Tomasulo's algorithm
\item \textbf{Speculative Execution}: Execute before branches resolved
\end{itemize}

\vspace{0.3cm}
\concept{Performance Benefits}:
\begin{itemize}
\item Higher instructions per cycle (IPC)
\item Better functional unit utilization
\item Reduced impact of individual hazards
\item Scalable performance improvements
\end{itemize}
\end{column}
\begin{column}{0.4\textwidth}
\begin{center}
\textbf{Superscalar Pipeline}
\vspace{0.3cm}

\tikz{
% Multiple issue pipeline
\node at (1.5,3.5) {\textbf{4-way Superscalar}};

% Fetch stage
\draw[thick] (0,3) rectangle (3,3.3);
\node at (1.5,3.15) {Fetch 4 Instructions};

% Decode stage
\foreach \x in {0, 0.75, 1.5, 2.25} {
  \draw[thick] (\x,2.5) rectangle (\x+0.6,2.8);
  \node at (\x+0.3,2.65) {Dec};
}

% Execute stage
\foreach \x/\unit in {0/ALU, 0.75/ALU, 1.5/LD/ST, 2.25/BR} {
  \draw[thick] (\x,2) rectangle (\x+0.6,2.3);
  \node at (\x+0.3,2.15) {\unit};
}

% Write back
\draw[thick] (0,1.5) rectangle (3,1.8);
\node at (1.5,1.65) {Write Back 4 Results};

% Arrows
\foreach \x in {0.3, 1.05, 1.8, 2.55} {
  \draw[->] (\x,3) -- (\x,2.8);
  \draw[->] (\x,2.5) -- (\x,2.3);
  \draw[->] (\x,2) -- (\x,1.8);
}

\node at (1.5,1.2) {\textbf{Up to 4 IPC}};
}
\end{center}
\end{column}
\end{columns}
\end{frame}

\begin{frame}{Out-of-Order Execution}
\begin{columns}
\begin{column}{0.6\textwidth}
\concept{Out-of-Order Execution Principle}:
\begin{itemize}
\item Instructions execute when operands ready
\item Not necessarily in program order
\item Improves functional unit utilization
\item Hides latency of slow operations
\end{itemize}

\vspace{0.3cm}
\concept{Key Components}:
\begin{itemize}
\item \textbf{Reservation Stations}: Hold instructions waiting for operands
\item \textbf{Reorder Buffer}: Maintain program order for commits
\item \textbf{Register Renaming}: Eliminate false dependencies
\item \textbf{Common Data Bus}: Broadcast results to waiting instructions
\end{itemize}

\vspace{0.3cm}
\concept{Tomasulo's Algorithm}:
\begin{itemize}
\item Dynamic scheduling with register renaming
\item Reservation stations track dependencies
\item Instructions execute when ready
\item Results committed in program order
\end{itemize}

\vspace{0.3cm}
\concept{Benefits and Costs}:
\begin{itemize}
\item \textbf{Benefits}: Higher IPC, latency tolerance
\item \textbf{Costs}: Complex hardware, power consumption
\item \textbf{Scalability}: Limited by window size
\item \textbf{Verification}: Complex correctness verification
\end{itemize}
\end{column}
\begin{column}{0.4\textwidth}
\begin{center}
\textbf{Out-of-Order Engine}
\vspace{0.3cm}

\tikz{
% Instruction flow
\node at (1.5,3.8) {\textbf{OoO Execution}};

% Fetch/Decode
\draw[thick] (0.5,3.4) rectangle (2.5,3.6);
\node at (1.5,3.5) {Fetch/Decode};

% Reservation stations
\draw[thick] (0,2.8) rectangle (1,3.2);
\node at (0.5,3) {RS1};
\draw[thick] (1.2,2.8) rectangle (2.2,3.2);
\node at (1.7,3) {RS2};
\draw[thick] (2.4,2.8) rectangle (3.4,3.2);
\node at (2.9,3) {RS3};

% Functional units
\draw[thick] (0,2.2) rectangle (1,2.6);
\node at (0.5,2.4) {ALU1};
\draw[thick] (1.2,2.2) rectangle (2.2,2.6);
\node at (1.7,2.4) {ALU2};
\draw[thick] (2.4,2.2) rectangle (3.4,2.6);
\node at (2.9,2.4) {LD/ST};

% Reorder buffer
\draw[thick] (0.5,1.6) rectangle (2.5,1.9);
\node at (1.5,1.75) {Reorder Buffer};

% Commit
\draw[thick] (0.5,1) rectangle (2.5,1.3);
\node at (1.5,1.15) {Commit (In-Order)};

% Arrows
\draw[->] (1.5,3.4) -- (0.5,3.2);
\draw[->] (1.5,3.4) -- (1.7,3.2);
\draw[->] (1.5,3.4) -- (2.9,3.2);
\draw[->] (0.5,2.8) -- (0.5,2.6);
\draw[->] (1.7,2.8) -- (1.7,2.6);
\draw[->] (2.9,2.8) -- (2.9,2.6);
\draw[->] (1.5,2.2) -- (1.5,1.9);
\draw[->] (1.5,1.6) -- (1.5,1.3);

\node at (1.7,0.7) {\textbf{Execute when ready}};
\node at (1.7,0.4) {\textbf{Commit in order}};
}
\end{center}
\end{column}
\end{columns}
\end{frame}

% subsection 6: Performance Analysis
\subsection{Performance Analysis}

\begin{frame}{ILP Performance Metrics}
\begin{columns}
\begin{column}{0.6\textwidth}
\concept{Key Performance Metrics}:
\begin{itemize}
\item \textbf{IPC (Instructions Per Cycle)}: Average instructions completed per cycle
\item \textbf{CPI (Cycles Per Instruction)}: Average cycles per instruction (1/IPC)
\item \textbf{Speedup}: Performance improvement over baseline
\item \textbf{Efficiency}: Actual performance vs. theoretical maximum
\end{itemize}

\vspace{0.3cm}
\concept{Factors Affecting ILP Performance}:
\begin{itemize}
\item \textbf{Instruction Mix}: Types of instructions in workload
\item \textbf{Dependency Frequency}: How often hazards occur
\item \textbf{Branch Frequency}: Control flow disruptions
\item \textbf{Cache Performance}: Memory system efficiency
\end{itemize}

\vspace{0.3cm}
\concept{Amdahl's Law for ILP}:
\begin{itemize}
\item Sequential portions limit parallel speedup
\item Dependencies create sequential bottlenecks
\item Diminishing returns with more parallel resources
\item Practical ILP limits: 2-4 instructions per cycle
\end{itemize}

\vspace{0.3cm}
\concept{Optimization Strategies}:
\begin{itemize}
\item Compiler optimizations for ILP
\item Hardware prediction improvements
\item Larger instruction windows
\item Better branch prediction
\end{itemize}
\end{column}
\begin{column}{0.4\textwidth}
\begin{center}
\textbf{ILP Performance}
\vspace{0.3cm}

% Performance graph
\tikz{
\draw[->] (0,0) -- (3,0);
\draw[->] (0,0) -- (0,3);
\node at (3.2,0) {Issue Width};
\node at (0,3.2) {IPC};

% Performance curve
\draw[thick, blue] (0.2,0.2) .. controls (1,1.5) and (2,2.2) .. (2.8,2.4);

% Theoretical maximum
\draw[thick, red, dashed] (0,0) -- (2.8,2.8);
\node at (2.2,2.5) {Theoretical};
\node at (2.2,2.2) {Actual};

% Practical limit
\draw[thick, green, dashed] (0,1.8) -- (3,1.8);
\node at (2.5,1.5) {Practical Limit};

% Labels
\node at (1.5,-0.3) {Superscalar Width};
\node at (-0.5,1.5) {IPC};
}

\vspace{0.5cm}

% Bottlenecks
\tikz{
\node at (1.5,1.5) {\textbf{ILP Bottlenecks}};
\node at (1.5,1.2) {• Dependencies};
\node at (1.5,0.9) {• Branch mispredictions};
\node at (1.5,0.6) {• Resource conflicts};
\node at (1.5,0.3) {• Memory latency};
}
\end{center}
\end{column}
\end{columns}
\end{frame}

% Summary
\subsection{Summary}

\begin{frame}{Chapter Summary}
\begin{columns}
\begin{column}{0.6\textwidth}
\concept{Key Concepts Covered}:
\begin{itemize}
\item Instruction Level Parallelism (ILP) fundamentals
\item Pipelining principles and implementation
\item Pipeline hazards and resolution techniques
\item Advanced ILP: superscalar and out-of-order execution
\item Performance analysis and optimization strategies
\item Trade-offs between complexity and performance
\end{itemize}

\vspace{0.3cm}
\concept{Technical Innovations}:
\begin{itemize}
\item Pipeline acceleration techniques
\item Data forwarding and hazard resolution
\item Branch prediction mechanisms
\item Dynamic instruction scheduling
\item Register renaming and speculation
\end{itemize}

\vspace{0.3cm}
\concept{Modern Relevance}:
\begin{itemize}
\item Foundation for all modern processors
\item Basis for high-performance computing
\item Principles in GPU and accelerator design
\item Continued evolution in processor architecture
\item Critical for energy-efficient computing
\end{itemize}
\end{column}
\begin{column}{0.4\textwidth}
\begin{center}
\textbf{ILP Evolution}
\vspace{0.3cm}

\tikz{
\node[draw, circle, minimum width=1.5cm] (center) at (0,0) {ILP};
\node[draw, rectangle] (pipe) at (0,1.5) {Pipelining};
\node[draw, rectangle] (super) at (-1.5,0.5) {Superscalar};
\node[draw, rectangle] (ooo) at (1.5,0.5) {Out-of-Order};
\node[draw, rectangle] (modern) at (0,-1.5) {Modern CPUs};

\draw[->] (pipe) -- (center);
\draw[->] (center) -- (super);
\draw[->] (center) -- (ooo);
\draw[->] (center) -- (modern);
}
\end{center}
\end{column}
\end{columns}
\end{frame}

\begin{frame}{Discussion Questions}
\begin{enumerate}
\item How does pipelining improve processor performance and what are its limitations?
\item What are the different types of pipeline hazards and how can they be resolved?
\item How does data forwarding reduce pipeline stalls and what are its costs?
\item What role does branch prediction play in maintaining pipeline efficiency?
\item How do superscalar processors achieve higher instruction throughput?
\item What are the benefits and challenges of out-of-order execution?
\item How do compiler optimizations complement hardware ILP techniques?
\item What factors limit the practical amount of ILP that can be extracted?
\end{enumerate}
\end{frame}


